% Zone „Das kennst du schon“ – aus bank/pythagoras/zone.jsonl (zusammenbau v0.1)
\einheitenkopf[zone][Kennst du schon]{Das kennst du schon}
\setcounter{aufgabe}{0}
%% TODO Grundvorstellungs-Aufgabe als erste Hauptnummer der Zone (2.8) fehlt in der Bank

%% TODO Titel: Ich-kann-Satz fehlt in der Bank (Platzhalter: Kettenname) – Nr. 1
%% TODO Anweisung: ein Satz über der ersten Teilaufgabe fehlt in der Bank – Nr. 1
\begin{aufgabe}{Quadrieren mit dem Taschenrechner, auch Dezimalzahlen; Quadratzahlen bis zwanzig hoch zwei aus dem Kopf; Quadrat vor Strich}
%% TODO Darstellung für schwach fehlt in der Bank (pythagoras-zone-f1-v1; Abschnitt „Für schwache Schüler“)
\swz{$11^2$ – wie viel?}{}
%% TODO Darstellung für schwach fehlt in der Bank (pythagoras-zone-f1-v3; Abschnitt „Für schwache Schüler“)
\swz{$1{,}4^2$ – wie viel?}{}
\end{aufgabe}

%% TODO Titel: Ich-kann-Satz fehlt in der Bank (Platzhalter: Kettenname) – Nr. 2
%% TODO Anweisung: ein Satz über der ersten Teilaufgabe fehlt in der Bank – Nr. 2
\begin{aufgabe}{Quadratwurzel mit dem Taschenrechner ziehen: erst die Summe oder Differenz unter der Wurzel ausrechnen, dann die Wurzel, dann den Näherungswert runden}
%% TODO Darstellung für schwach fehlt in der Bank (pythagoras-zone-f2-v1; Abschnitt „Für schwache Schüler“)
\swz{$\sqrt{121}$ – wie viel?}{}
%% TODO Darstellung für schwach fehlt in der Bank (pythagoras-zone-f2-v3; Abschnitt „Für schwache Schüler“)
\swz{$\sqrt{50}$ – auf eine Stelle gerundet?}{}
\end{aufgabe}

%% TODO Titel: Ich-kann-Satz fehlt in der Bank (Platzhalter: Kettenname) – Nr. 3
%% TODO Anweisung: ein Satz über der ersten Teilaufgabe fehlt in der Bank – Nr. 3
\begin{aufgabe}{Rechtwinkliges Dreieck erkennen: Rechtwinkelmarke (Bogen mit Punkt) in der Skizze finden, Dreiecksarten}
\begin{teile}
\teil Rechtwinklig? \janein
\end{teile}
%% TODO Darstellung für schwach fehlt in der Bank (pythagoras-zone-f3-v3; Abschnitt „Für schwache Schüler“)
\swz{Winkel $35^\circ$ und $55^\circ$ – dritter Winkel? Rechtwinklig?}{}
\end{aufgabe}

%% TODO Titel: Ich-kann-Satz fehlt in der Bank (Platzhalter: Kettenname) – Nr. 4
%% TODO Anweisung: ein Satz über der ersten Teilaufgabe fehlt in der Bank – Nr. 4
\begin{aufgabe}{Längeneinheiten: Ergebnis mit Einheit, cm und m vor dem Rechnen angleichen, Runden auf eine Dezimale}
%% TODO Darstellung für schwach fehlt in der Bank (pythagoras-zone-f4-v1; Abschnitt „Für schwache Schüler“)
\swz{$3$ m – wie viel cm?}{}
%% TODO Darstellung für schwach fehlt in der Bank (pythagoras-zone-f4-v3; Abschnitt „Für schwache Schüler“)
\swz{$6{,}847$ m – auf eine Dezimale?}{}
\end{aufgabe}

%% TODO Titel: Ich-kann-Satz fehlt in der Bank (Platzhalter: Kettenname) – Nr. 5
%% TODO Anweisung: ein Satz über der ersten Teilaufgabe fehlt in der Bank – Nr. 5
\begin{aufgabe}{Formel nach einer Größe umstellen mit Strich, hier c² = a² + b² nach a²}
%% TODO Darstellung für schwach fehlt in der Bank (pythagoras-zone-f5-v1; Abschnitt „Für schwache Schüler“)
\swz{$g = e + f$ – nach $e$ umstellen?}{}
%% TODO Darstellung für schwach fehlt in der Bank (pythagoras-zone-f5-v3; Abschnitt „Für schwache Schüler“)
\swz{$P = Q + R$ mit $P = 7{,}5$ und $R = 2{,}8$ – $Q$?}{}
\end{aufgabe}

%% TODO Titel: Ich-kann-Satz fehlt in der Bank (Platzhalter: Kettenname) – Nr. 6
%% TODO Anweisung: ein Satz über der ersten Teilaufgabe fehlt in der Bank – Nr. 6
\begin{aufgabe}{Koordinaten ablesen und Differenzen bilden, auch mit negativen Zahlen}
\swz{Koordinaten von $P$?}{\begin{ksys}[xmin=-1,xmax=5,ymin=-1,ymax=5,ablesen] \punkt{3}{4}{P} \end{ksys}}
%% TODO Darstellung für schwach fehlt in der Bank (pythagoras-zone-f6-v3; Abschnitt „Für schwache Schüler“)
\swz{x-Werte $-3$ und $4$ – Unterschied?}{}
\end{aufgabe}

%% TODO Titel: Ich-kann-Satz fehlt in der Bank (Platzhalter: Kettenname) – Nr. 7
%% TODO Anweisung: ein Satz über der ersten Teilaufgabe fehlt in der Bank – Nr. 7
\begin{aufgabe}{Höhe, Seitenhöhe, Seitenkante, Mantellinie und Radius an Pyramide und Kegel benennen, Radius aus dem Durchmesser}
%% TODO Darstellung für schwach fehlt in der Bank (pythagoras-zone-f7-v1; Abschnitt „Für schwache Schüler“)
\swz{Durchmesser $14$ cm – Radius?}{}
%% TODO Darstellung für schwach fehlt in der Bank (pythagoras-zone-f7-v3; Abschnitt „Für schwache Schüler“)
\swz{Durchmesser $9{,}6$ cm – Radius?}{}
\end{aufgabe}

%% TODO Titel: Ich-kann-Satz fehlt in der Bank (Platzhalter: Kettenname) – Nr. 8
%% TODO Anweisung: ein Satz über der ersten Teilaufgabe fehlt in der Bank – Nr. 8
\begin{aufgabe}{Eigenschaften von gleichschenkligem Dreieck, gleichschenkligem Trapez, Parallelogramm, Drachen und Rechteck: Höhe, Symmetrieachse, Diagonalen}
\begin{teile}
\teil Rechteck: Diagonalen gleich lang? \janein
\end{teile}
%% TODO Darstellung für schwach fehlt in der Bank (pythagoras-zone-f8-v3; Abschnitt „Für schwache Schüler“)
\swz{Gleichschenkliges Dreieck, Grundseite $11$ cm – Abstand Höhenfußpunkt bis Ecke?}{}
\end{aufgabe}

%% TODO Titel: Ich-kann-Satz fehlt in der Bank (Platzhalter: Kettenname) – Nr. 9
%% TODO Anweisung: ein Satz über der ersten Teilaufgabe fehlt in der Bank – Nr. 9
\begin{aufgabe}{Strecke halbieren, halbe Differenz zweier Seiten bilden (Überstand am Trapez)}
%% TODO Darstellung für schwach fehlt in der Bank (pythagoras-zone-f9-v1; Abschnitt „Für schwache Schüler“)
\swz{$18$ cm – die Hälfte?}{}
%% TODO Darstellung für schwach fehlt in der Bank (pythagoras-zone-f9-v3; Abschnitt „Für schwache Schüler“)
\swz{$9{,}4$ m und $4{,}6$ m – halbe Differenz?}{}
\end{aufgabe}

%% TODO Titel: Ich-kann-Satz fehlt in der Bank (Platzhalter: Kettenname) – Nr. 10
%% TODO Anweisung: ein Satz über der ersten Teilaufgabe fehlt in der Bank – Nr. 10
\begin{aufgabe}{Quadratwurzel mit dem Taschenrechner ziehen: erst die Summe oder Differenz unter der Wurzel ausrechnen, dann die Wurzel, dann den Näherungswert runden – Fehler finden}
\swz[3]{Ole rechnet: \rechnung{\sqrt{121 + 3\,600} &= 11 + 60 \\ &= 71} – wo ist der Fehler?}{}
\end{aufgabe}

%% TODO Titel: Ich-kann-Satz fehlt in der Bank (Platzhalter: Kettenname) – Nr. 11
%% TODO Anweisung: ein Satz über der ersten Teilaufgabe fehlt in der Bank – Nr. 11
\begin{aufgabe}{Quadratwurzel mit dem Taschenrechner ziehen: erst die Summe oder Differenz unter der Wurzel ausrechnen, dann die Wurzel, dann den Näherungswert runden – selbst rechnen}
%% TODO Darstellung für schwach fehlt in der Bank (pythagoras-zone-f2-v6; Abschnitt „Für schwache Schüler“)
\swz{$\sqrt{784 + 2\,025}$ – wie viel?}{}
\end{aufgabe}
